\documentclass[11pt, a4paper]{article}

\usepackage{amsmath}
\usepackage{amssymb}
\usepackage{geometry}
\usepackage{graphicx}
\usepackage{hyperref}
\usepackage{listings}
\usepackage{xcolor}

\geometry{margin=1in}

\definecolor{codegreen}{rgb}{0,0.6,0}
\definecolor{codegray}{rgb}{0.5,0.5,0.5}
\definecolor{codepurple}{rgb}{0.58,0,0.82}
\definecolor{backcolour}{rgb}{0.95,0.95,0.92}

\lstdefinestyle{mystyle}{
    backgroundcolor=\color{backcolour},   
    commentstyle=\color{codegreen},
    keywordstyle=\color{magenta},
    numberstyle=\tiny\color{codegray},
    stringstyle=\color{codepurple},
    basicstyle=\ttfamily\footnotesize,
    breakatwhitespace=false,         
    breaklines=true,                 
    captionpos=b,                    
    keepspaces=true,                 
    numbers=left,                    
    numbersep=5pt,                  
    showspaces=false,                
    showstringspaces=false,
    showtabs=false,                  
    tabsize=2
}
\lstset{style=mystyle}

\title{\textbf{Solid Mechanics Numerical Implementation Report} \\ \Large A Comparative Analysis of Theoretical Physics and Numerical Methods}
\author{Solid Mechanics Simulator Engine}
\date{\today}

\begin{document}

\maketitle

\begin{abstract}
This report outlines the core solid mechanics physics derived from the lecture syllabus (Stress/Strain Tensors, Beam Bending, Axisymmetric Problems, Section Properties, and Column Buckling) and maps them directly to the numerical implementations utilized in the Python FastAPI and PyQt6 math engine.
\end{abstract}

\tableofcontents
\newpage

\section{Introduction}
The objective of the interactive solid mechanics laboratory is to bridge analytical solid mechanics formulations with discrete numerical evaluation. The application handles multi-dimensional tensors, differential integrations for deflection, and failure criterion evaluations using \texttt{numpy} arrays and linear algebra subroutines.

\section{Module 1: Stress, Strain, and Tensors}
\subsection{Principal Stresses and Invariants (Lectures 7-10)}
\textbf{Theory:} Principal stresses ($\sigma_1, \sigma_2, \sigma_3$) are found by solving the eigenvalue problem for the Cauchy stress tensor $\boldsymbol{\sigma}$:
\begin{equation}
\det(\boldsymbol{\sigma} - \lambda \mathbf{I}) = 0
\end{equation}
Where the roots $\lambda_i$ are the principal stresses.

\textbf{Numerical Implementation:} 
Instead of symbolically solving the cubic characteristic polynomial, the engine utilizes LAPACK routines via \texttt{numpy.linalg.eig()}, which numerically computes eigenvalues via QR decomposition.
\begin{lstlisting}[language=Python]
# tensor_math.py
eigenvalues, eigenvectors = np.linalg.eig(tensor)
idx = eigenvalues.argsort()[::-1]   
ps = eigenvalues[idx] # [sigma_1, sigma_2, sigma_3]
\end{lstlisting}

\subsection{Arbitrary Traction Vectors (Lecture 2)}
\textbf{Theory:} The traction vector $\mathbf{T}$ on a plane defined by normal $\mathbf{\hat{n}}$ is given by Cauchy's stress theorem:
\begin{equation}
\mathbf{T} = \boldsymbol{\sigma} \mathbf{\hat{n}}
\end{equation}
Normal stress $\sigma_n = \mathbf{T} \cdot \mathbf{\hat{n}}$ and shear stress $\tau_n = \sqrt{|\mathbf{T}|^2 - \sigma_n^2}$.

\textbf{Numerical Implementation:}
\begin{lstlisting}[language=Python]
n_vec = np.array(n_vec)
n_hat = n_vec / np.linalg.norm(n_vec)
T = np.dot(tensor, n_hat)
sigma_n = np.dot(T, n_hat)
tau_n = np.sqrt(np.linalg.norm(T)**2 - sigma_n**2)
\end{lstlisting}

\section{Module 2: Beam Bending and Deflection}
\subsection{Euler-Bernoulli Deflection (Lecture 27)}
\textbf{Theory:} Deflection $v(x)$ is defined by integrating the bending moment $M(x)$ twice:
\begin{equation}
\frac{d^2v}{dx^2} = \frac{M(x)}{EI} \implies v(x) = \iint \frac{M(x)}{EI} dx dx + C_1 x + C_2
\end{equation}

\textbf{Numerical Implementation:}
The beam is discretized into $N$ points. Integration is handled via cumulative trapezoidal summation. Constants $C_1$ and $C_2$ are explicitly solved using the geometric boundary conditions (e.g., $v(0)=0$ and $v(L)=0$ for simply supported beams).
\begin{lstlisting}[language=Python]
import scipy.integrate as integrate
theta = integrate.cumtrapz(M / (E * I), x, initial=0)
v = integrate.cumtrapz(theta, x, initial=0)
# Enforce boundary conditions mathematically
v_corrected = v - (v[-1]/L) * x 
\end{lstlisting}

\section{Module 3: Axisymmetric Problems}
\subsection{Thick-Walled Pressure Vessels (Lecture 20)}
\textbf{Theory:} Lam\'{e}'s equations for internal pressure $p_i$ yield radial and hoop stresses:
\begin{equation}
\sigma_r = \frac{p_i r_i^2}{r_o^2 - r_i^2} \left(1 - \frac{r_o^2}{r^2}\right), \quad \sigma_\theta = \frac{p_i r_i^2}{r_o^2 - r_i^2} \left(1 + \frac{r_o^2}{r^2}\right)
\end{equation}

\textbf{Numerical Implementation:}
Evaluated over an array representing the radial continuum $r \in [r_i, r_o]$.
\begin{lstlisting}[language=Python]
r_vals = np.linspace(ri, ro, 50)
term1 = (p * ri**2) / (ro**2 - ri**2)
sigma_r = term1 * (1 - (ro**2 / r_vals**2))
sigma_t = term1 * (1 + (ro**2 / r_vals**2))
\end{lstlisting}

\section{Module 4: Section Properties \& Shear Center}
\subsection{Geometric Exactness (Lecture 26)}
\textbf{Theory:} The exact moment of inertia of a composite shape is defined by the parallel axis theorem $I_x = \sum (I_{xi} + A_i d_i^2)$. The shear center eccentricity $e$ for a thin-walled C-channel is $e = \frac{b^2 h^2 t}{4 I_x}$.

\textbf{Numerical Implementation:}
The math engine avoids thin-walled approximations for fundamental geometry by computing the distinct rectangular sub-components.
\begin{lstlisting}[language=Python]
# Exact Centroid
Qy = A_web * (t/2) + 2 * (A_flange * (b/2))
xc = Qy / A_total

# Exact Inertia using Parallel Axis
Ix_flange = (b * t**3)/12.0 + A_flange * d_flange**2
Ix_total = Ix_web + 2 * Ix_flange
\end{lstlisting}

\section{Module 5: Column Buckling}
\subsection{Euler and Johnson Criteria (Lecture 28)}
\textbf{Theory:} Slender columns buckle at $P_{cr} = \frac{\pi^2 EI}{(KL)^2}$. Intermediate columns yield according to the Johnson Parabola: $P_{cr} = A \left[ S_y - \frac{S_y^2 (KL/r)^2}{4 \pi^2 E} \right]$.

\textbf{Numerical Implementation:}
Conditional execution branching maps the critical slenderness ratio against the material's yield limits.
\begin{lstlisting}[language=Python]
slenderness = (K * L) / r
sig_cr_euler = (math.pi**2 * E) / (slenderness**2)

if sig_cr_euler > Sy / 2:
    sig_cr = Sy - (Sy**2 * slenderness**2) / (4 * math.pi**2 * E)
else:
    sig_cr = sig_cr_euler
\end{lstlisting}

\section{Module 6: Theories of Failure}
\subsection{Von Mises \& Tresca Criteria (Lectures 31-32)}
\textbf{Theory:} 
Tresca: $\sigma_{failure} = \sigma_1 - \sigma_3$ \\
Von Mises: $\sigma_{VM} = \sqrt{\frac{1}{2} \left[ (\sigma_1 - \sigma_2)^2 + (\sigma_2 - \sigma_3)^2 + (\sigma_3 - \sigma_1)^2 \right]}$

\textbf{Numerical Implementation:}
These are executed immediately following the \texttt{np.linalg.eig} routine in the principal stress domain to evaluate the live Factor of Safety.
\begin{lstlisting}[language=Python]
tresca = ps[0] - ps[2]
vm = np.sqrt(0.5 * ((ps[0]-ps[1])**2 + (ps[1]-ps[2])**2 + (ps[2]-ps[0])**2))
FS_tresca = Sy / tresca
FS_vm = Sy / vm
\end{lstlisting}

\end{document}
